\documentclass[10pt,a4paper]{amsart}
\usepackage[margin=19mm]{geometry}
\usepackage{amsmath,amssymb,mathtools}
\usepackage[scr=rsfs,cal=euler]{mathalfa}
\usepackage{xfrac}
\usepackage[unicode,hidelinks,bookmarks=false]{hyperref}
\newcommand{\ie}{i.e.\ }
\newcommand{\cf}{cf.\ }
\newcommand{\resp}{respectively\ }
\newcommand{\Subsection}{\subsection}
\providecommand{\nobreakdash}{\nobreak\hbox{-}}
\setlength{\emergencystretch}{2em}
\multlinegap0pt
\usepackage{bm}
\usepackage{bbm}
\usepackage{dsfont}
\usepackage{xspace}
\usepackage{cancel}
\usepackage{mathtools}
\usepackage{amsthm}
\makeatletter
\@ifundefined{theorem}{\newtheorem{theorem}{Theorem}}{}
\makeatother
\title[Balanced Fibonacci word rectangles, and beyond]{Balanced Fibonacci word rectangles, and beyond}
\author{Jeffrey Shallit, Ingrid Vukusic}
\date{Source: arXiv:2509.25994v6 (CC BY 4.0)}
\begin{document}
\maketitle

\noindent\emph{Source and licence.} Balanced Fibonacci word rectangles, and beyond. Version arXiv:2509.25994v6 (CC BY 4.0), distributed under the Creative Commons Attribution 4.0 International licence, as recorded in the arXiv licence metadata for this exact version and shown on the abstract page https://arxiv.org/abs/2509.25994v6. Full rights notice: licenses/arxiv-2509.25994-CC-BY-4.0.txt.\par\smallskip

\noindent\textit{Editorial scope.} Counted unit: the Theorem whose source label is \texttt{None} in the article (source0.tex lines 682--685, source label \texttt{None}), delivered together with the article's own complete original proof of it (source0.tex lines 687--756, one \texttt{proof} environment of 70 lines). No mathematics is written, repaired or re-derived: statement and proof are the article's own text line for line. Required context retained uncounted: none. Every definition, display and label of those lines is retained unchanged. Omitted from this package and not redistributed: 1--681, 686--686, 757--840. Nothing inside a retained excerpt is omitted or altered: every retained body is delivered exactly as the article's lines are written, so no omission receipt of any kind accompanies this package. Citations: the retained excerpts carry no citation command at all, so no bibliography entry of the article is transcribed and no reference is redistributed. Reference adaptation: none was needed and none was made; every reference inside the delivered unit resolves inside it, so no unresolved reference or citation is produced. Printed slips kept literal: no printed word, symbol, hypothesis or number of the counted statement or of its proof was altered. Layout and class: the article's own class is \texttt{dmtcs-episciences} with packages including inputenc, subfigure, natbib, amsthm, mathtools; the wrapper is the lane's pinned \texttt{amsart} editorial wrapper at 10pt on A4, which restates the theorem-like environments the retained text needs and adds the article's own macro definitions that the retained text needs (none), with extra wrapper packages (bm, bbm, dsfont, xspace, cancel, mathtools). The delivered numbering is the unit's own. Assets: no retained span contains a figure environment, a graphics-inclusion command, a PDF-page inclusion, a table or a TikZ picture; no image file is copied into this package, the package declares no asset dependency and no image file is redistributed with it. Licence: version arXiv:2509.25994v6 of this article is distributed under the Creative Commons Attribution 4.0 International licence, as recorded in the arXiv licence metadata for this exact version and shown on the abstract page https://arxiv.org/abs/2509.25994v6. A full rights notice is delivered at licenses/arxiv-2509.25994-CC-BY-4.0.txt. Characters: every retained excerpt is pure ASCII, so the delivered engine, XeLaTeX with its default Latin Modern text font, reports no missing character.
\begin{theorem}
    There is an automaton of $92$ states that on input $m$ and $n$ in base $2$ determines the balance of the $m \times n$ rectangles
    built from $\bf t$.
\end{theorem}

\begin{proof}
We use the following {\tt Walnut} code.
The idea is illustrated for the case $i,m,n$ all even.  We consider pairing all the entries of the matrix with an even index with the odd index that follows it.
This leaves unpaired entries in the leftmost column and the
rightmost column.

In the leftmost column these entries are
$t_{i+1} + t_{i+3} + \cdots + t_{i+m-1}$.
Since all the indices are odd,
this is the same as
$(1-t_{i/2}) +  (1-t_{(i+2)/2}) + \cdots + (1-t_{(i+m-2)/2})
= m/2 - (t_{i/2} + t_{i/2+1} + \cdots + t_{(i+m)/2-1})$.
In the rightmost column these entries are
$t_{i+n} + t_{i+n+2} + \cdots + t_{i+m+n-2} =
t_{(i+n)/2} + t_{(i+n)/2 + 1} + \cdots +
t_{(i+m+n)/2 -1}$.
The case of $i$ odd is similar, and we reduce the other cases to this one by removing a row or column, or both.

Here the automaton {\tt TM2} computes the Thue-Morse sequence, but with input given in base-$(-2)$.  The automaton
{\tt twonegtwo} takes two inputs,
$w$ and $x$ in parallel, and accepts if $w$ considered as a base-$(-2)$ expansion equals $x$ considered as a base-$2$ expansion.  Both can be downloaded from the website of the first author.
\begin{verbatim}
def n2even "?msd_neg_2 Ek n=2*k":
def n2odd "?msd_neg_2 Ek n=2*k+1":
def prefsum "?msd_neg_2 (s=n/2 & $n2even(n)) | 
   (s=n/2+TM2[n-1] & $n2odd(n))":
# sum of t[0..n-1]
# 17 states
def factorsum "?msd_neg_2 Eu,v $prefsum(i,u) & $prefsum(i+n,v) & 
   s+u=v":
# sum of t[i..i+n-1]
# 317 states
def m_even_n_even "?msd_neg_2 Ea,b $factorsum(i/2,m/2,a) & 
   $factorsum((i+n)/2,m/2,b) & s=2*(b-a)":
# s = 2T(i,m,n) - mn  for m even, n even

def m_even_n_odd "?msd_neg_2 Ea,b $m_even_n_even(i,m,n-1,a) & 
   $factorsum(i+n-1,m,b) & s=a+2*b-m":
def m_odd_n_even "?msd_neg_2 Ea,b $m_even_n_even(i,m-1,n,a) & 
   $factorsum(i+m-1,n,b) & s=a+2*b-n":
def m_odd_n_odd "?msd_neg_2 Ea,b $m_odd_n_even(i,m,n-1,a) & 
   $factorsum(i+n-1,m,b) & s=a+2*b-m":

def tmt "?msd_neg_2  i>=0 & m>=0 & n>=0 & 
(($n2even(m) & $n2even(n) & $m_even_n_even(i,m,n,s)) |
($n2even(m) & $n2odd(n) & $m_even_n_odd(i,m,n,s)) |
($n2odd(m) & $n2even(n) & $m_odd_n_even(i,m,n,s)) |
($n2odd(m) & $n2odd(n) & $m_odd_n_odd(i,m,n,s)))":
# computes 2T(i,m,n) - mn for all i, m, n >= 0
# 3678 states

def maxval "?msd_neg_2 (Ei $tmt(i,m,n,x)) & 
   Ai,s $tmt(i,m,n,s) => s<=x":
# 81 states

def tmbal "Ex,y,b $twonegtwo((?msd_neg_2 x),m) & 
   $twonegtwo((?msd_neg_2 y),n) & $twonegtwo((?msd_neg_2 b),z) &
   $maxval(?msd_neg_2 x,y,b)":
# 94 states

def aut0 "$tmbal(m,n,0)":
def aut1 "$tmbal(m,n,1)":
def aut2 "$tmbal(m,n,2)": 
def aut3 "$tmbal(m,n,3)":
def aut4 "$tmbal(m,n,4)":

combine TMB aut0=0 aut1=1 aut2=2 aut3=3 aut4=4:
# result has 92 states 
\end{verbatim}\vspace{-\baselineskip}
\end{proof}

\end{document}
